\chapter{Problem Formulation}
\label{chap:problem}

This chapter provides the mathematical formulation of the joint user pairing and power allocation problem in downlink PD-NOMA systems.

\section{System Overview}

Consider a single-cell downlink PD-NOMA system with:
\begin{itemize}
    \item One base station (BS) with total transmit power $P$
    \item $N$ single-antenna User Equipment (UE) indexed by $\mathcal{N} = \{1, 2, \ldots, N\}$
    \item $M$ Physical Resource Blocks (PRBs) with bandwidth $B$ each
    \item Perfect Channel State Information (CSI) at the BS
    \item Quasi-static flat fading channels
\end{itemize}

\section{Channel Model}

The channel coefficient between the BS and user $i$ is:
\begin{equation}
h_i = \sqrt{PL_i \cdot SF_i} \cdot g_i
\label{eq:channel_gain}
\end{equation}

where:
\begin{itemize}
    \item $PL_i$ is the path loss (distance-dependent attenuation)
    \item $SF_i$ is the log-normal shadow fading
    \item $g_i \sim \mathcal{CN}(0, 1)$ is the small-scale Rayleigh fading coefficient
\end{itemize}

The channel power gain is:
\begin{equation}
|h_i|^2 = PL_i \cdot SF_i \cdot |g_i|^2
\end{equation}

We assume users are ordered such that $|h_1|^2 \leq |h_2|^2 \leq \cdots \leq |h_N|^2$ (weak to strong).

\section{NOMA Signal Model}

\subsection{Two-User Pairing}

For simplicity, we focus on two-user clusters. Consider pairing user $i$ (weak) with user $j$ (strong) where $|h_i|^2 < |h_j|^2$ on a single PRB with power $P$.

The transmitted signal is:
\begin{equation}
s = \sqrt{\alpha P} \cdot x_i + \sqrt{(1-\alpha) P} \cdot x_j
\end{equation}

where:
\begin{itemize}
    \item $x_i, x_j$ are unit-power symbols for users $i, j$
    \item $\alpha \in (0, 1)$ is the power allocation coefficient
    \item $\alpha > 0.5$ allocates more power to the weaker user $i$
\end{itemize}

\subsection{User $i$ (Weak User) Reception}

User $i$ treats user $j$'s signal as noise. The received signal is:
\begin{equation}
y_i = h_i \left( \sqrt{\alpha P} x_i + \sqrt{(1-\alpha) P} x_j \right) + n_i
\end{equation}

where $n_i \sim \mathcal{CN}(0, \sigma^2)$ is AWGN with power $\sigma^2$.

The Signal-to-Interference-plus-Noise Ratio (SINR) is:
\begin{equation}
\text{SINR}_i = \frac{\alpha P |h_i|^2}{(1-\alpha) P |h_i|^2 + \sigma^2}
\label{eq:sinr_weak}
\end{equation}

\subsection{User $j$ (Strong User) Reception}

User $j$ first decodes user $i$'s signal treating its own as noise:
\begin{equation}
\text{SINR}_{j \to i} = \frac{\alpha P |h_j|^2}{(1-\alpha) P |h_j|^2 + \sigma^2}
\label{eq:sinr_sic}
\end{equation}

After successfully decoding $x_i$, user $j$ performs SIC to cancel $x_i$ from the received signal, then decodes its own:
\begin{equation}
\text{SINR}_j = \frac{(1-\alpha) P |h_j|^2}{\sigma^2}
\label{eq:sinr_strong}
\end{equation}

\subsection{SIC Feasibility Constraint}

For successful SIC, user $j$ must decode $x_i$ with sufficient reliability:
\begin{equation}
\text{SINR}_{j \to i} \geq \gamma_{\text{th}}
\label{eq:sic_constraint}
\end{equation}

where $\gamma_{\text{th}}$ is the SIC decoding threshold (typically 8 dB = 6.31 linear).

Substituting Eq.~\eqref{eq:sinr_sic}:
\begin{equation}
\frac{\alpha P |h_j|^2}{(1-\alpha) P |h_j|^2 + \sigma^2} \geq \gamma_{\text{th}}
\end{equation}

Rearranging:
\begin{equation}
\alpha \geq \frac{\gamma_{\text{th}} \left((1-\alpha) P |h_j|^2 + \sigma^2\right)}{P |h_j|^2}
\end{equation}

For high SNR ($P |h_j|^2 \gg \sigma^2$), this simplifies to:
\begin{equation}
\frac{|h_j|^2}{|h_i|^2} \geq \frac{\gamma_{\text{th}}}{1 - \gamma_{\text{th}}/(1 + \gamma_{\text{th}})} \approx \gamma_{\text{th}}
\label{eq:channel_disparity}
\end{equation}

Thus, SIC requires sufficient channel gain disparity (typically $|h_j|^2 / |h_i|^2 \geq 8$ dB).

\section{Rate Analysis}

\subsection{Individual Rates}

Using Shannon capacity formula with bandwidth $B$:

\textbf{Weak user $i$ rate:}
\begin{equation}
R_i = B \log_2\left(1 + \frac{\alpha P |h_i|^2}{(1-\alpha) P |h_i|^2 + \sigma^2}\right)
\label{eq:rate_weak}
\end{equation}

\textbf{Strong user $j$ rate} (after SIC):
\begin{equation}
R_j = B \log_2\left(1 + \frac{(1-\alpha) P |h_j|^2}{\sigma^2}\right)
\label{eq:rate_strong}
\end{equation}

\subsection{Sum-Rate}

The sum-rate for pair $(i, j)$ is:
\begin{equation}
R_{\text{sum}}(i, j, \alpha) = R_i + R_j
\label{eq:sum_rate}
\end{equation}

\section{Power Allocation Optimization}

For a fixed pair $(i, j)$, the optimal power allocation maximizes sum-rate:
\begin{equation}
\begin{aligned}
\alpha^* = \argmax_{\alpha \in (0,1)} \quad & R_{\text{sum}}(i, j, \alpha) \\
\text{s.t.} \quad & \text{SINR}_{j \to i}(\alpha) \geq \gamma_{\text{th}}
\end{aligned}
\label{eq:power_opt}
\end{equation}

\subsection{Closed-Form Solution}

Taking the derivative of $R_{\text{sum}}$ with respect to $\alpha$ and setting to zero yields a transcendental equation. In practice, we use:

\begin{enumerate}
    \item \textbf{Bisection Search:} Numerically find $\alpha^*$ in range $[\alpha_{\min}, 1)$ where $\alpha_{\min}$ satisfies the SIC constraint.
    
    \item \textbf{Fixed Allocation:} Common heuristic sets $\alpha = 0.8$ for all pairs.
    
    \item \textbf{Water-Filling:} Allocates power proportional to channel quality difference.
\end{enumerate}

For this project, we use bisection search with tolerance $10^{-4}$ and maximum 80 iterations.

\section{User Pairing Problem}

\subsection{Decision Variables}

Define binary pairing variables:
\begin{equation}
x_{ij} = \begin{cases}
1 & \text{if users } i \text{ and } j \text{ are paired} \\
0 & \text{otherwise}
\end{cases}
\end{equation}

For symmetric pairing, $x_{ij} = x_{ji}$ and $i \neq j$.

\subsection{Optimization Formulation}

The joint pairing and power allocation problem is:
\begin{equation}
\begin{aligned}
\max_{x, \alpha} \quad & \sum_{i=1}^{N} \sum_{j=i+1}^{N} x_{ij} \cdot R_{\text{sum}}(i, j, \alpha_{ij}) \\
\text{s.t.} \quad & \sum_{j=1, j \neq i}^{N} x_{ij} \leq 1, \quad \forall i \in \mathcal{N} \quad \text{(matching constraint)} \\
& x_{ij} \in \{0, 1\}, \quad \forall i, j \in \mathcal{N} \\
& \frac{|h_j|^2}{|h_i|^2} \geq \gamma_{\text{th}}, \quad \text{if } x_{ij} = 1 \quad \text{(SIC constraint)} \\
& |\theta_i - \theta_j| \geq \Delta\theta_{\min}, \quad \text{if } x_{ij} = 1 \quad \text{(angular separation)} \\
& \alpha_{ij} \in (0, 1), \quad \forall i, j
\end{aligned}
\label{eq:main_problem}
\end{equation}

where $\theta_i$ is the angular position of user $i$.

\subsection{Complexity Analysis}

Problem~\eqref{eq:main_problem} is a Mixed-Integer Nonlinear Program (MINLP):
\begin{itemize}
    \item Binary variables $x_{ij}$: $\binom{N}{2}$ decisions
    \item Continuous variables $\alpha_{ij}$: require solving \eqref{eq:power_opt} for each pair
    \item Matching constraint makes it a maximum-weight matching problem
\end{itemize}

\textbf{Computational Complexity:}
\begin{itemize}
    \item Exhaustive search: $O((N-1)!! \cdot N) = O(N!)$ for all possible matchings
    \item Blossom algorithm: $O(N^3)$ for maximum-weight matching
    \item Heuristic methods: $O(N \log N)$ for sorting-based approaches
\end{itemize}

For real-time scheduling (1 ms TTI in 5G NR), $N=100$ users require solutions in <1 ms, making optimization infeasible.

\section{Fairness Metrics}

\subsection{Jain's Fairness Index}

Measures rate distribution uniformity:
\begin{equation}
\mathcal{J} = \frac{\left(\sum_{i=1}^{N} R_i\right)^2}{N \sum_{i=1}^{N} R_i^2}
\label{eq:jains_index}
\end{equation}

$\mathcal{J} = 1$ indicates perfect fairness; $\mathcal{J} \to 1/N$ indicates one user monopolizes resources.

\subsection{Proportional Fairness}

Maximizes the sum of logarithmic rates:
\begin{equation}
\max \sum_{i=1}^{N} \log(R_i)
\end{equation}

This balances throughput and fairness, preventing starvation of weak users.

\section{Reformulation as Graph Learning}

\subsection{Graph Representation}

We model the pairing problem as an undirected graph $\mathcal{G} = (\mathcal{V}, \mathcal{E})$:

\begin{itemize}
    \item \textbf{Nodes} $\mathcal{V}$: Each user $i$ is a node with feature vector:
    \begin{equation}
    \mathbf{x}_i = [|h_i|^2, \, d_i, \, \theta_i, \, \text{SNR}_i, \, \text{PL}_i]^\top
    \end{equation}
    
    \item \textbf{Edges} $\mathcal{E}$: Edge $(i, j)$ exists if:
    \begin{enumerate}
        \item $|h_j|^2 / |h_i|^2 \geq \gamma_{\text{th}}$ (SIC feasibility)
        \item $|\theta_i - \theta_j| \geq \Delta\theta_{\min}$ (angular separation)
        \item $i < j$ (to avoid duplicates)
    \end{enumerate}
\end{itemize}

\subsection{Learning Objectives}

For each edge $(i, j) \in \mathcal{E}$, we predict:

\begin{enumerate}
    \item \textbf{Pairing Probability} $p_{ij} \in [0, 1]$: Likelihood that $(i, j)$ should be paired
    
    \item \textbf{Sum-Rate} $\hat{R}_{ij}$: Predicted achievable sum-rate
    
    \item \textbf{Power Allocation} $\hat{\alpha}_{ij} \in (0, 1)$: Predicted optimal power split
\end{enumerate}

This multi-task formulation enables joint optimization and exploits coupling between objectives.

\section{Summary}

This chapter formulated the NOMA user pairing problem as:
\begin{itemize}
    \item A constrained MINLP with exponential search space
    \item A maximum-weight matching problem solvable in $O(N^3)$ via Blossom
    \item A graph representation learning problem amenable to GNNs
\end{itemize}

The next chapters detail system implementation, channel modeling, and solution methodologies.
